\documentclass[10pt,a4paper]{amsart}
\usepackage[margin=19mm]{geometry}
\usepackage{amsmath,amssymb,mathtools}
\usepackage[scr=rsfs,cal=euler]{mathalfa}
\usepackage{xfrac}
\usepackage[unicode,hidelinks,bookmarks=false]{hyperref}
\newcommand{\ie}{i.e.\ }
\newcommand{\cf}{cf.\ }
\newcommand{\resp}{respectively\ }
\newcommand{\Subsection}{\subsection}
\providecommand{\nobreakdash}{\nobreak\hbox{-}}
\setlength{\emergencystretch}{2em}
\multlinegap0pt
\usepackage{amssymb}
\newtheorem{lemma}{Lemma}
\newtheorem{theorem}{Theorem}
\newtheorem{proposition}{Proposition}
\def\Om{\Omega}
\def\bZ{\mathbb{Z}}
\def\cG{\mathcal{G}}
\def\de{\delta}
\def\si{\sigma}
\def\xiGE{\xi_{\mathrm{Mu}}}
\def\xiTM{\xi_{\mathrm{TM}}}
\def\xipd{\xi_{\mathrm{pd}}}
\title{Thue-Morse sequence and groups of intermediate growth}
\author{Rostislav Grigorchuk, Yaroslav Vorobets}
\date{Source: arXiv:2605.30605v1 (CC BY 4.0)}
\begin{document}
\maketitle

\noindent\emph{Source and licence.} Thue-Morse sequence and groups of intermediate growth. Version arXiv:2605.30605v1 (CC BY 4.0), distributed by arXiv under the Creative Commons Attribution 4.0 International licence, http://creativecommons.org/licenses/by/4.0/, as recorded in the arXiv licence metadata for this exact version and shown on the abstract page https://arxiv.org/abs/2605.30605v1. Full licence notice: \texttt{licenses/arxiv-2605.30605-CC-BY-4.0.txt}.\par\smallskip

\noindent\textit{Editorial scope.} Counted unit: the article's proposition two-more-inj (source0.tex lines 2710--2727). The article's complete original proof of it is delivered with it (source0.tex lines 2729--2825, one \texttt{proof} environment of 97 lines). No mathematics is written, repaired or re-derived: the statement and the proof are the article's own lines, in the article's own notation, and no retained line is substituted or re-flowed.  Retained uncounted as the source-local context the proof needs: lines 727--731; lines 1031--1034; lines 1082--1085; lines 1390--1394; lines 1527--1531; lines 1550--1556; lines 1588--1593; lines 1895--1900.  Nothing was omitted from any retained span, so the package carries no omission record.  No retained line contains a citation, so the wrapper carries no bibliography and no bibliography entry.  No retained line contains a figure environment, an image-inclusion command or drawing code, so no image is delivered and the package declares no asset dependency.  Editorial formatting only, in addition: an AMS \texttt{amsart} preamble that restates the article's own declarations and macros used by the retained lines, namely the environments \texttt{lemma}, \texttt{theorem}, \texttt{proposition} on separate counters, together with the standard packages those lines need.  The retained environments print in this document as Lemma 1, Theorem 1, Proposition 1 in order of appearance, whereas the article prints the same items with the numbering of its own full document.  The complete native source of the article as distributed in the arXiv e-print for this exact version is bundled unchanged as \texttt{sources/201684/source0.tex}.
\begin{lemma}\label{xi-GE}
Let $n=2^km$, where $k$ is a nonnegative integer and $m$ is a positive odd 
integer.  Then the $n$-th letter of the string $\xiGE$ is $a$ if $k=0$, $c$ 
if $k$ is odd, and $d$ if $k$ is even and positive.
\end{lemma}

\begin{lemma}\label{block-factor-1}
The $2$-block factor map $f_{\phi_1}$ is two-to-one when restricted to the 
set $\Om(\xiTM)$, and $f_{\phi_1}(\Om(\xiTM))=\Om(\xipd)$.
\end{lemma}

\begin{lemma}\label{block-factor-2}
The $1$-block factor map $f_{\phi_2}$ is one-to-one when restricted to the 
set $\Om(\xiGE)$, and $f_{\phi_2}(\Om(\xiGE))=\Om(\xipd)$.
\end{lemma}

\begin{theorem}\label{TFG-embeds}
Suppose that $T_1$ and $T_2$ are aperiodic homeomorphisms of Cantor sets 
and $T_2$ is a continuous factor of $T_1$.  Then the topological full group 
$[[T_2]]$ embeds into $[[T_1]]$.
\end{theorem}

\begin{lemma}\label{model-isomorphism}
Suppose that the homeomorphism $T$ is aperiodic.  Then the homomorphism 
$H_T$ is one-to-one.  Moreover, the restriction of $H_T$ to the group 
$C^\times_T(X,\bZ)$ is an isomorphism onto the group $[[T]]$.
\end{lemma}

\begin{lemma}\label{monoid-embed}
Suppose $T_1:X_1\to X_1$ and $T_2:X_2\to X_2$ are homeomorphisms of Cantor 
sets and $f:X_1\to X_2$ is a continuous function such that $fT_1=T_2f$.  
Then the map $\Phi_f$ is a homomorphism of the monoid $C(X_2,\bZ)$ with 
operation $\oplus_{T_2}$ to the monoid $C(X_1,\bZ)$ with operation 
$\oplus_{T_1}$.
\end{lemma}

\begin{lemma}\label{group-embed}
Under the assumptions of Lemma \ref{monoid-embed}, the restriction of 
$\Phi_f$ to the group $C^\times_{T_2}(X_2,\bZ)$ is a homomorphism to the 
group $C^\times_{T_1}(X_1,\bZ)$.  This group homomorphism is one-to-one 
provided that the function $f$ is onto.
\end{lemma}

\begin{proposition}\label{inj-homomorph}
There exists a one-to-one group homomorphism 
$h:\cG_{(01)^\infty}\to[[\si|\Omega(\xiGE)]]$ such that $h(a)=\de_a$, 
$h(c)=\de_c$ and $h(d)=\de_d$, where $c=c_{(01)^\infty}$, 
$d=d_{(01)^\infty}$.
\end{proposition}

\begin{proposition}\label{two-more-inj}
(i) There exists a one-to-one group homomorphism $h_1:\cG_{(01)^\infty}\to 
[[\si_{\{x,y\}}|\Om(\xipd)]]$ such that 
$h_1(z)=\de_{U_z\cap\Om(\xipd);S_1}$ for all $z\in\{a,c,d\}$, where 
$S_1=\si_{\{x,y\}}|\Om(\xipd)$, $U_a=[.xy]_{\{x,y\}}\cup [.xxxy]_{\{x,y\}}$,
$U_c=[.y]_{\{x,y\}}$ and $U_d=[.xxy]_{\{x,y\}}$.

(ii) There exists a one-to-one group homomorphism $h_2:\cG_{(01)^\infty}\to 
[[\si_{\{0,1\}}|\Om(\xiTM)]]$ such that 
$h_2(z)=\de_{W_z\cap\Om(\xiTM);S_2}$ for all $z\in\{a,c,d\}$, where 
$S_2=\si_{\{0,1\}}|\Om(\xiTM)$,
\begin{eqnarray*}
W_a &=& [.011]_{\{0,1\}}\cup [.100]_{\{0,1\}}\cup [.01011]_{\{0,1\}}\cup 
[.10100]_{\{0,1\}},\\
W_c &=& [.00]_{\{0,1\}}\cup [.11]_{\{0,1\}},\\
W_d &=& [.0100]_{\{0,1\}}\cup [.1011]_{\{0,1\}}.
\end{eqnarray*}
\end{proposition}

\begin{proof}
Suppose $T_1:X_1\to X_1$ and $T_2:X_2\to X_2$ are aperiodic homeomorphisms 
of Cantor sets, and $f:X_1\to X_2$ is a continuous map that is onto and 
satisfies $fT_1=T_2f$.  By Theorem \ref{TFG-embeds}, there exists a 
one-to-one homomorphism $\Psi:[[T_2]]\to[[T_1]]$.  An explicit description 
of $\Psi$ can be extracted from Lemmas \ref{model-isomorphism} and 
\ref{group-embed}.  Namely, if $F\in[[T_2]]$ is a homeomorphism given by 
$F(x)=T_2^{\nu(x)}(x)$, $x\in X_2$, where $\nu:X_2\to\bZ$ is a continuous 
function, then its image $\Psi(F)$ is given by 
$\bigl(\Psi(F)\bigr)(x)=T_1^{\nu(f(x))}(x)$ for all $x\in X_1$.  In the 
case when $F=\de_{U;T_2}$ for some clopen set $U\subset X_2$, we obtain 
that $\Psi(F)=\de_{f^{-1}(U);T_1}$.

By Proposition \ref{inj-homomorph}, there is a one-to-one homomorphism 
$h:\cG_{(01)^\infty}\to[[\si_{\{a,c,d\}}|\Om(\xiGE)]]$ such that 
$h(z)=\de_z$ for all $z\in\{a,c,d\}$.  Recall that $\de_z$ is the short 
notation for the generalized $2$-cycle 
$\de_{[.z]_{\{a,c,d\}}\cap\Om(\xiGE);S}$, where 
$S=\si_{\{a,c,d\}}|\Om(\xiGE)$.  Let 
$f_{\phi_2}:\{a,c,d\}^{\bZ}\to\{x,y\}^{\bZ}$ be the block factor map 
considered in Lemma \ref{block-factor-2}.  In view of the lemma, the 
restriction of $f_{\phi_2}$ to the set $\Om(\xiGE)$ can be regarded a 
homeomorphism onto $\Om(\xipd)$.  Let $g_1$ be the inverse of that 
homeomorphism.  Then $g_1S_1=Sg_1$, where $S_1=\si_{\{x,y\}}|\Om(\xipd)$.
Hence $g_1$ induces a one-to-one homomorphism $\Psi_1:[[S]]\to[[S_1]]$ as 
described above.  Then $h_1=\Psi_1h$ is a one-to-one homomorphism of the 
group $\cG_{(01)^\infty}$ to $[[S_1]]$.  For any $z\in\{a,c,d\}$ we obtain 
that $h_1(z)=\Psi_1(\de_z)=\de_{U'_z;S_1}$, where $U'_z= 
g_1^{-1}\bigl([.z]_{\{a,c,d\}}\cap\Om(\xiGE)\bigr)
=f_{\phi_2}\bigl([.z]_{\{a,c,d\}}\cap\Om(\xiGE)\bigr)$.

It follows from the description of the string $\xiGE$ (Lemma \ref{xi-GE}) 
that any occurrence of the letter $d$ in $\xiGE$ is followed by $ac$.  
Hence $[.d]_{\{a,c,d\}}\cap\Om(\xiGE)=[.dac]_{\{a,c,d\}}\cap\Om(\xiGE)$.  
Any occurrence of the letter $a$ is followed by either $c$ or $dac$.  Hence 
the set $[.a]_{\{a,c,d\}}\cap\Om(\xiGE)$ is the union of 
$[.ac]_{\{a,c,d\}}\cap\Om(\xiGE)$ and $[.adac]_{\{a,c,d\}}\cap\Om(\xiGE)$.  
Recall that the $1$-block factor map $f_{\phi_2}$ is induced by the 
function $\phi_2:\{a,c,d\}\to\{x,y\}$ given by $\phi_2(a)=\phi_2(d)=x$ and 
$\phi_2(c)=y$.  Therefore $f_{\phi_2}\bigl([.c]_{\{a,c,d\}}\bigr)\subset 
[.y]_{\{x,y\}}$, $f_{\phi_2}\bigl([.ac]_{\{a,c,d\}}\bigr)\subset 
[.xy]_{\{x,y\}}$, $f_{\phi_2}\bigl([.dac]_{\{a,c,d\}}\bigr)\subset 
[.xxy]_{\{x,y\}}$ and $f_{\phi_2}\bigl([.adac]_{\{a,c,d\}}\bigr)\subset 
[.xxxy]_{\{x,y\}}$.  Note that $[.c]_{\{a,c,d\}}$, $[.ac]_{\{a,c,d\}}$, 
$[.dac]_{\{a,c,d\}}$ and $[.adac]_{\{a,c,d\}}$ are disjoint sets that cover 
$\Om(\xiGE)$.  As observed in the proof of Lemma \ref{block-factor-2}, the 
word $xxxx$ is not a subword of the string $\xipd$.  Hence the set 
$[.xxxx]_{\{x,y\}}$ is disjoint from $\Om(\xipd)$.  Consequently, 
$\Om(\xipd)$ is covered by disjoint sets $[.y]_{\{x,y\}}$, 
$[.xy]_{\{x,y\}}$, $[.xxy]_{\{x,y\}}$ and $[.xxxy]_{\{x,y\}}$.  Since the 
restriction of $f_{\phi_2}$ to $\Om(\xiGE)$ is a one-to-one map onto 
$\Om(\xipd)$, it follows that 
$f_{\phi_2}\bigl([.c]_{\{a,c,d\}}\cap\Om(\xiGE)\bigr)=
[.y]_{\{x,y\}}\cap\Om(\xipd)$, 
$f_{\phi_2}\bigl([.ac]_{\{a,c,d\}}\cap\Om(\xiGE)\bigr)= 
[.xy]_{\{x,y\}}\cap\Om(\xipd)$, 
$f_{\phi_2}\bigl([.dac]_{\{a,c,d\}}\cap\Om(\xiGE)\bigr)=
[.xxy]_{\{x,y\}}\cap\Om(\xipd)$ and 
$f_{\phi_2}\bigl([.adac]_{\{a,c,d\}}\cap\Om(\xiGE)\bigr)=
[.xxxy]_{\{x,y\}}\cap\Om(\xipd)$.  As a consequence, for any 
$z\in\{a,c,d\}$ we have 
$f_{\phi_2}\bigl([.z]_{\{a,c,d\}}\cap\Om(\xiGE)\bigr)=U_z\cap\Om(\xipd)$, 
where $U_z$ is as defined in the proposition.  This completes the proof of 
the statement (i).

We proceed to the statement (ii).  Let 
$f_{\phi_1}:\{0,1\}^{\bZ}\to\{x,y\}^{\bZ}$ be the block factor map 
considered in Lemma \ref{block-factor-1}.  The map $f_{\phi_1}$ is 
two-to-one when restricted to $\Om(\xiTM)$ and $f_{\phi_1}(\Om(\xiTM))= 
\Om(\xipd)$.  It follows from the proof of Lemma \ref{block-factor-1} that 
$f_{\phi_1}$ is two-to-one on $\{0,1\}^{\bZ}$.  As a consequence, 
$f_{\phi_1}^{-1}(\Om(\xipd))=\Om(\xiTM)$.  Let $g_2$ be the restriction of 
$f_{\phi_1}$ to $\Om(\xiTM)$.  Then $g_2S_2=S_1g_2$, where 
$S_2=\si_{\{0,1\}}|\Om(\xiTM)$.  Hence $g_2$ induces a one-to-one 
homomorphism $\Psi_2:[[S_1]]\to[[S_2]]$ as described above.  Then 
$h_2=\Psi_2h_1$ is a one-to-one homomorphism of the group 
$\cG_{(01)^\infty}$ to $[[S_2]]$.  For any $z\in\{a,c,d\}$ we have
$h_2(z)=\Psi_2\bigl(\de_{U_z\cap\Om(\xipd);S_1}\bigr)
=\de_{g_2^{-1}(U_z\cap\Om(\xipd));S_2}$.  Note that 
$g_2^{-1}(U_z\cap\Om(\xipd))=f_{\phi_1}^{-1}(U_z)\cap\Om(\xiTM)$ since 
$f_{\phi_1}^{-1}(\Om(\xipd))=\Om(\xiTM)$.  Recall that the $2$-block factor
map $f_{\phi_1}$ is induced by the function $\phi_1:\{0,1\}^2\to\{x,y\}$
given by $\phi_1(0,1)=\phi_1(1,0)=x$ and 
$\phi_1(0,0)=\phi_1(1,1)=y$.  We obtain that
\begin{align*}
&f_{\phi_1}^{-1}([.y]_{\{x,y\}}) = [.00]_{\{0,1\}}\cup [.11]_{\{0,1\}},\\
&f_{\phi_1}^{-1}([.xy]_{\{x,y\}}) = [.011]_{\{0,1\}}\cup 
[.100]_{\{0,1\}},\\
&f_{\phi_1}^{-1}([.xxy]_{\{x,y\}}) = [.0100]_{\{0,1\}}\cup
[.1011]_{\{0,1\}},\\
&f_{\phi_1}^{-1}([.xxxy]_{\{x,y\}}) = [.01011]_{\{0,1\}}\cup 
[.10100]_{\{0,1\}}.
\end{align*}
It follows that for any $z\in\{a,c,d\}$ we have $f_{\phi_1}^{-1}(U_z)=W_z$, 
where $W_z$ is as defined in the proposition.  This completes the proof of 
the statement (ii).
\end{proof}

\end{document}
